\section{Three Objects as One: Unitary Equivalence}
\begin{quote}\small
\noindent\textbf{Status.} \emph{Legacy structural-identity block inside the active main body.}\par
\noindent\textbf{Block function.} This section compresses the unitary-equivalence reading that ties the three-object architecture into one common structural carrier.
\end{quote}
\label{sec:unitary}
%% ============================================================

\begin{definition}[Sector classification from biquaternion]
The three sectors correspond to the sign of $\mathrm{Re}(\zminus)$:
\begin{itemize}
  \item \textbf{Matter} ($K_1$): $\mathrm{Re}(z_1)<0$, boost $\beta<0$, spacelike.
  \item \textbf{Tachyon} ($K_2=J\cdot K_1$): $\mathrm{Re}(z_2)>0$, boost $\beta>0$, timelike.
  \item \textbf{Photon} ($K_3=e_1 e_2$): $\mathrm{Re}(z_3)=0$, boost $\beta=0$, lightlike.
\end{itemize}
\end{definition}

\begin{theorem}[Unitary equivalence via Stone--von Neumann]\label{thm:unitary-equiv}
The three kernels $K_1,K_2,K_3$ carry unitarily equivalent
irreducible representations of the Weyl algebra $\WA$.
The boost $J$ implements the equivalence between $K_1$ and $K_2$,
and the compensator $C$ produces $K_3$ from $K_1$.
\end{theorem}

\begin{proof}
By Stone--von Neumann~\cite{brendecke2026rh}, every irreducible
representation of $\WA$ is unitarily equivalent to $L^2(\mathbb{Q}^+\backslash A_1)$.
All three kernels are subspaces of $\Hplus$ which carries this
representation.
Since $J$ is unitary and maps $K_1\leftrightarrow K_2$, they are
equivalent.
$K_3=C(K_1)$ lies in $\Hplusplus$, the unique $J$-invariant subspace,
which is again irreducible (no proper $J$-invariant subspace exists by
T4+T5 of~\cite{brendecke2026rh}).
\end{proof}

%% ============================================================

